% =============================================================================
% Fixtures/sample-ieeetran-manuscript.tex
%
% FICTIONAL SAMPLE MANUSCRIPT written for the ScholarOS test suite.
% Every system, dataset, experiment, number, person, affiliation and reference
% named in this file is invented. Nothing here reports real research, and no
% text was copied from any real manuscript.
% The .bib file named at the end is intentionally not shipped; `outline` does not need it.
% Purpose: an IEEEtran-style input of realistic size and structure for
% `scholaros outline` (Scripts/acceptance/m0-acceptance.sh, row 1). It contains
% numbered and starred \section / \subsection headings only, so the heading
% count can be cross-checked with
%     grep -c -E '\\(section|subsection)\*?\{' Fixtures/sample-ieeetran-manuscript.tex
% =============================================================================
\documentclass[conference]{IEEEtran}

\usepackage{amsmath}
\usepackage{amssymb}
\usepackage{graphicx}
\usepackage{booktabs}
\usepackage{cite}

\newcommand{\sysname}{\textsc{Lanternfish}}
\newcommand{\dataset}{\textsc{TidePool-9}}
\newcommand{\ms}{\,\mathrm{ms}}
\newcommand{\kb}{\,\mathrm{kB}}

\begin{document}

\title{Lanternfish: Scheduling Tidal Sensor Buoys With Drifting Clocks\\
{\large A Fictional Sample Manuscript for Parser Tests}}

\author{\IEEEauthorblockN{Mira Okonkwo-Lindqvist}
\IEEEauthorblockA{Department of Imaginary Systems\\
Saltmarsh Institute of Technology\\
Port Averill, Nowhere\\
mira@saltmarsh.example}
\and
\IEEEauthorblockN{Teodor Vashti}
\IEEEauthorblockA{Coastal Computing Laboratory\\
University of Greywater\\
Greywater, Nowhere\\
teodor@greywater.example}
\and
\IEEEauthorblockN{Anselm Ruiz-Harcourt}
\IEEEauthorblockA{Buoy Networks Group\\
Saltmarsh Institute of Technology\\
Port Averill, Nowhere\\
anselm@saltmarsh.example}}

\maketitle

\begin{abstract}
Low-power sensor buoys wake on a schedule, sample the water column, and go back
to sleep. When their clocks drift apart, two buoys that should take turns on a
shared radio channel start talking over each other. This invented paper
describes \sysname, a scheduler that tolerates clock drift without a global time
source. \sysname\ lets each buoy keep a small table of its neighbours' observed
wake times and shifts its own slot by a bounded amount each cycle. On the
made-up \dataset\ trace we report fewer collisions and a longer battery life
than two imaginary baselines. All results in this document are fabricated for
the purpose of exercising a LaTeX parser.
\end{abstract}

\begin{IEEEkeywords}
sensor networks, clock drift, duty cycling, scheduling, fictional example
\end{IEEEkeywords}

\section{Introduction}
\label{sec:intro}

A tidal buoy spends most of its life asleep. Each wake-up costs energy, and each
radio transmission costs more. Operators therefore assign every buoy a short
slot in a repeating cycle, and the buoys are expected to stay inside their
slots for months at a time~\cite{okonkwo2031slots}.

The assumption behind slotted schedules is that all buoys agree on what time it
is. In practice they do not. A crystal oscillator on a cold buoy runs slower
than the same part on a warm one, and the difference accumulates. After a few
weeks, two neighbouring buoys may wake at the same moment and corrupt each
other's packets~\cite{vashti2030drift,harcourt2029coldwater}.

\subsection{Motivating Scenario}
\label{sec:intro-scenario}

Consider twelve buoys moored along an estuary. The upstream buoys sit in fresh,
cold water, while the downstream ones are warmed by a shallow mudflat every
afternoon. Their clocks diverge by a few seconds per week. A shore station can
correct them, but only two buoys are within range of it.

\subsection{Why Existing Approaches Fall Short}
\label{sec:intro-gap}

A buoy could carry a satellite receiver and read the time directly, but the
receiver draws more current than the rest of the board combined. Alternatively
the shore station could flood time beacons through the network, but beacons are
themselves radio traffic and collide like any other packet~\cite{ruiz2030beacon}.

\subsection{Contributions}
\label{sec:intro-contrib}

This fictional paper makes three contributions.
\begin{itemize}
  \item A neighbour table that records when each nearby buoy was last heard,
        using only packets the buoy would have received anyway.
  \item A slot-shifting rule with a provable bound on how far a buoy may move
        in one cycle (Section~\ref{sec:design-shift}).
  \item An evaluation on the invented \dataset\ trace
        (Section~\ref{sec:eval}).
\end{itemize}

\subsection{Paper Organization}
\label{sec:intro-org}

Section~\ref{sec:background} reviews duty cycling and clock drift.
Section~\ref{sec:design} presents the design of \sysname.
Section~\ref{sec:impl} describes the prototype.
Section~\ref{sec:eval} reports the evaluation, Section~\ref{sec:discussion}
discusses limitations, and Section~\ref{sec:conclusion} concludes.

\section{Background and Related Work}
\label{sec:background}

\subsection{Duty-Cycled Radios}
\label{sec:bg-duty}

A duty-cycled radio is on for a fraction $\delta$ of each cycle of length $T$.
The energy spent per cycle is approximately
\begin{equation}
  E_{\mathrm{cycle}} = \delta T P_{\mathrm{on}} + (1 - \delta) T P_{\mathrm{sleep}},
  \label{eq:energy}
\end{equation}
where $P_{\mathrm{on}}$ and $P_{\mathrm{sleep}}$ are the power draw in each
state. Shrinking $\delta$ saves energy only as long as the buoys still manage
to hear one another.

\subsection{Clock Drift in Marine Deployments}
\label{sec:bg-drift}

Drift is usually quoted in parts per million. A drift of $40$\,ppm amounts to
roughly $3.5$ seconds per day. Temperature is the dominant cause in the
deployments imagined here~\cite{harcourt2029coldwater}, followed by supply
voltage sag as the battery ages.

\subsection{Reference-Broadcast Synchronization}
\label{sec:bg-rbs}

Reference-broadcast schemes have a designated node transmit a beacon, and all
receivers timestamp its arrival. The receivers then compare timestamps among
themselves. The approach removes sender-side uncertainty but still requires
extra packets~\cite{ruiz2030beacon,pemberton2028pairwise}.

\subsection{Schedule Repair Without Synchronization}
\label{sec:bg-repair}

A different line of invented work abandons synchronization altogether and
repairs the schedule after collisions are detected~\cite{abara2031repair}.
\sysname\ sits between the two camps: it never computes a shared clock, yet it
acts before collisions occur.

\subsection*{A Note on Terminology}

We use \emph{slot} for the interval in which one buoy may transmit, \emph{cycle}
for one full rotation through all slots, and \emph{guard} for the idle time
between adjacent slots.

\section{System Design}
\label{sec:design}

\subsection{Design Goals}
\label{sec:design-goals}

\sysname\ was designed around four goals.
\begin{enumerate}
  \item \textbf{No extra packets.} Timing information must be derived from
        traffic that already exists.
  \item \textbf{Bounded movement.} A buoy may not shift its slot by more than a
        fixed fraction of the guard interval per cycle.
  \item \textbf{Local state only.} Each buoy stores information about its
        one-hop neighbours and nothing else.
  \item \textbf{Graceful failure.} If a neighbour disappears, the remaining
        buoys must not chase its last known position forever.
\end{enumerate}

\subsection{Neighbour Table}
\label{sec:design-table}

Every buoy keeps a table with one row per neighbour. A row stores the
neighbour's identifier, the local time at which its most recent packet started,
and a confidence counter that decays when the neighbour is silent. The table
occupies $11$ bytes per row in our imagined implementation.

\subsection{Slot Shifting Rule}
\label{sec:design-shift}

Let $s_i$ be the start of buoy $i$'s slot in its own local time, and let
$p_i$ and $n_i$ be the observed start times of the slots immediately before and
after it. The buoy computes the midpoint
\begin{equation}
  m_i = \frac{p_i + n_i}{2}
  \label{eq:midpoint}
\end{equation}
and moves towards it by a bounded step:
\begin{equation}
  s_i \leftarrow s_i + \operatorname{clip}\bigl(\alpha (m_i - s_i),\, -g/4,\, g/4\bigr),
  \label{eq:shift}
\end{equation}
where $g$ is the guard interval and $0 < \alpha \le 1$ is a damping factor.

\subsection{Convergence Argument}
\label{sec:design-converge}

If every buoy applies (\ref{eq:shift}) once per cycle and drift per cycle is
smaller than $g/4$, the spacing between adjacent slots cannot shrink below
$g/2$. A sketch of the argument is given in the appendix. The full proof is
omitted because this manuscript is fictional.

\subsection{Handling Departures and Arrivals}
\label{sec:design-churn}

When the confidence counter of a neighbour reaches zero, its row is deleted and
the buoy treats the next surviving neighbour as adjacent. A newly deployed buoy
listens for two full cycles, picks the widest gap it observed, and starts in
the middle of that gap.

\subsection{Interaction With the Shore Station}
\label{sec:design-shore}

The shore station is treated as an ordinary neighbour whose slot never moves.
Buoys within range of it are thereby anchored, and the anchoring propagates
outward one hop per cycle.

\section{Implementation}
\label{sec:impl}

\subsection{Hardware Platform}
\label{sec:impl-hw}

The imagined prototype runs on a buoy controller with a $16$\,MHz
microcontroller, $32\kb$ of RAM, and a sub-gigahertz radio. Timekeeping uses a
$32.768$\,kHz watch crystal with a stated tolerance of $\pm 20$\,ppm at room
temperature.

\subsection{Firmware Structure}
\label{sec:impl-fw}

The scheduler is a single interrupt handler of about three hundred lines. It
runs at the end of each receive window, updates the neighbour table, and
programs the wake-up timer for the next cycle. Table~\ref{tab:footprint}
summarizes the made-up memory footprint.

\begin{table}[t]
  \caption{Firmware footprint of the fictional prototype.}
  \label{tab:footprint}
  \centering
  \begin{tabular}{lrr}
    \toprule
    Component & Flash (kB) & RAM (bytes) \\
    \midrule
    Radio driver & 6.2 & 148 \\
    Neighbour table & 0.9 & 176 \\
    Slot scheduler & 1.4 & 52 \\
    Sensor sampling & 3.8 & 96 \\
    \midrule
    Total & 12.3 & 472 \\
    \bottomrule
  \end{tabular}
\end{table}

\subsection{Timestamp Capture}
\label{sec:impl-ts}

Packet start times are captured by the radio's start-of-frame interrupt, which
latches the timer value in hardware. Software latency therefore does not enter
the measurement.

\subsection*{Build and Flash Procedure}

The firmware is built with a standard cross-compiler and written to the buoy
over a two-wire debug port. No step in the procedure requires network access.

\section{Evaluation}
\label{sec:eval}

\subsection{Experimental Setup}
\label{sec:eval-setup}

We replay the invented \dataset\ trace, which describes twelve buoys over
ninety simulated days. Water temperature at each mooring follows a daily
sinusoid with a mooring-specific offset. Each configuration is run with ten
random seeds.

\subsection{Baselines}
\label{sec:eval-baselines}

We compare against two imaginary baselines. \textsc{Fixed} assigns slots once
and never changes them. \textsc{Beacon} floods a time beacon from the shore
station every hour and resets all clocks to it.

\subsection{Collision Rate}
\label{sec:eval-collisions}

Table~\ref{tab:collisions} lists the fraction of slots in which two or more
buoys transmitted at once. \textsc{Fixed} degrades steadily as drift
accumulates. \textsc{Beacon} stays low but pays for it in beacon traffic.
\sysname\ remains below both after the first week.

\begin{table}[t]
  \caption{Collision rate on the fictional \dataset\ trace (per cent of slots).}
  \label{tab:collisions}
  \centering
  \begin{tabular}{lrrr}
    \toprule
    Day range & \textsc{Fixed} & \textsc{Beacon} & \sysname \\
    \midrule
    1--7 & 0.4 & 0.6 & 0.7 \\
    8--30 & 3.9 & 0.8 & 0.5 \\
    31--60 & 11.2 & 0.9 & 0.5 \\
    61--90 & 19.6 & 1.1 & 0.6 \\
    \bottomrule
  \end{tabular}
\end{table}

\subsection{Energy Consumption}
\label{sec:eval-energy}

Using (\ref{eq:energy}) with the duty cycles observed in simulation, the
projected battery life is $212$ days for \textsc{Fixed}, $171$ days for
\textsc{Beacon}, and $205$ days for \sysname. The gap between \textsc{Fixed}
and \sysname\ comes from the slightly longer receive window that \sysname\
needs in order to hear both neighbours.

\subsection{Convergence After Deployment}
\label{sec:eval-converge}

A buoy added on day $45$ settles into a stable slot after a median of six
cycles. The slowest of the ten seeds needed fourteen cycles.

\subsection{Sensitivity to the Damping Factor}
\label{sec:eval-alpha}

We sweep $\alpha$ from $0.1$ to $1.0$. Small values converge slowly; values
above $0.8$ oscillate when two adjacent buoys correct towards each other in
the same cycle. We use $\alpha = 0.5$ elsewhere in this section.

\subsection{Scalability}
\label{sec:eval-scale}

Repeating the experiment with $48$ and $96$ simulated buoys leaves the
collision rate within $0.2$ percentage points of the twelve-buoy result, because
each buoy only ever reasons about two neighbours.

\subsection*{Threats to Validity}

Every number above comes from a simulation of an invented trace. Real moorings
experience fouling, storms, and vandalism, none of which are modelled.

\section{Discussion}
\label{sec:discussion}

\subsection{Limitations}
\label{sec:disc-limits}

The slot-shifting rule assumes that a buoy can hear both of its schedule
neighbours. In a sparse line topology the slot before and the slot after may
belong to buoys that are out of radio range, in which case the rule has nothing
to work with.

\subsection{Security Considerations}
\label{sec:disc-security}

A malicious transmitter could pull its neighbours' slots by sending packets at
misleading times. Authenticating packets prevents impersonation but not replay
of a genuine packet at a later moment.

\subsection{Deployment Lessons}
\label{sec:disc-lessons}

Because the deployment is imaginary, the lessons are too. We would expect the
confidence counter to need tuning per site, since a buoy that is submerged at
high tide is silent for hours without having failed.

\subsection{Future Work}
\label{sec:disc-future}

We plan to extend the neighbour table to two hops, to study schedules with
unequal slot lengths, and to test the damping rule on drifting ice floes, which
do not exist in this paper's universe either.

\section{Conclusion}
\label{sec:conclusion}

This fictional paper presented \sysname, a scheduler that keeps duty-cycled
buoys from colliding as their clocks drift, using only timing information taken
from ordinary traffic. On an invented trace it kept collisions under one per
cent for ninety days at a small cost in receive time.

\section*{Acknowledgment}

The authors thank the nonexistent harbour master of Port Averill for access to
moorings that were never installed.

\section*{Data Availability}

There is no data. The \dataset\ trace is a name used in this sample manuscript
and does not correspond to any file.

\appendices

\section{Sketch of the Spacing Bound}
\label{app:proof}

\subsection{Setup}
\label{app:proof-setup}

Order the buoys by slot start and write $d_i = s_{i+1} - s_i$ for the spacing
between consecutive slots. Initially $d_i = g$ for all $i$.

\subsection{One-Cycle Step}
\label{app:proof-step}

In one cycle each of $s_i$ and $s_{i+1}$ moves by at most $g/4$ because of the
clip in (\ref{eq:shift}). If drift contributed nothing, $d_i$ could therefore
shrink by at most $g/2$. The midpoint rule in (\ref{eq:midpoint}) moves a buoy
away from whichever neighbour is closer, so two adjacent buoys never both move
towards each other by the full step unless their outer neighbours are far away,
in which case the outer gaps absorb the movement in the next cycle.

\section{Simulation Parameters}
\label{app:params}

\subsection{Radio Model}
\label{app:params-radio}

Packets are lost independently with probability $0.02$. Two packets that
overlap in time by more than one symbol are both lost.

\subsection{Temperature Model}
\label{app:params-temp}

Mooring $k$ has temperature
$\theta_k(t) = \bar{\theta}_k + A_k \sin(2 \pi t / 24\,\mathrm{h})$,
with $\bar{\theta}_k$ between $4$ and $17$ degrees and $A_k$ between $0.5$ and
$6$ degrees. Drift is a quadratic function of the distance from $25$ degrees.

\bibliographystyle{IEEEtran}
\bibliography{sample-ieeetran-references}

\end{document}
